\input ctustyle3
\worktype [O/EN]
\faculty {F8}
\department {Department of Theoretical Computer Science}
\workname{Semester project for the BI-TEX course}
\title {Document Background Grid \nl for Measuring Typesetting in Various Units}
\author {Adam Barla}
\date {SS 2020/2021}
\abstractEN {This document describes use and development of \OpTeX{} macro. This macro creates a grid on the background of each page of the document. The size and placement of the grid can be defined by the user. Grid keeps the Typesetting of the original document. It is designed to help with measuring a relative and the absolute positon of texts, pictures, charts and other elements on the page.}
\abstractSK {Tento dokument popisuje použitie a vývoj makra pre \OpTeX{}. Toto makro vytvorí na pozadí každej strany dokumentu mriežku ľubovoľnej velkosti s ľubobolným umiestnením. Mriežka nezasahuje do pôvodnej sadzby dokumentu. Pomocou tejto mriežky je jednoduché odčítať vzájomnú ale aj absolutnú polohu všetkých textov, obrázkov, tabuliek a iných elementov na stránke.}
% \declaration {I declare that I did not slack off.}
\makefront

\chap Introduction
The \TeX{} program is used in academic environments primarily thanks to the quality of its mathematical typesetting. However, it excels in many other areas as well. One of them, for example, is the line-breaking algorithm. What contributes to this excellence is that \TeX{} processes the entire source text at once. The author therefore only sees the final typesetting after the program finishes running. When fine-tuning the final appearance of the text, it is necessary to repeatedly re-run the program and look for possible errors. This can initially be overwhelming for a new user, especially if they are accustomed to text editors such as MS Word, where changes are visible immediately. This work describes a macro for \OpTeX{} designed to help find causes of incorrect typesetting placement or to determine the exact position of text in a document.

\chap Manual
This chapter serves as a guide for achieving the desired results when using the macro.
The grid itself consists of vertical and horizontal lines whose spacing can be defined using any unit in \TeX{}. On the top and left sides of the grid, every tenth line is numbered. Numbered lines are also extended and thicker for convenient distance reading. Regardless of thickness, the distance from the grid's origin is read on the left or top side of the vertical or horizontal line, respectively.
\sec Default Usage
After downloading the macro "grid.tex", simply include it in your document and activate it. You can achieve this using these two commands:
\begtt
\input grid.tex
\grid ,(,,,)
\endtt
The result will be a grid covering the entire page except for the margins. Only the numbering of every tenth line on both the vertical and horizontal axes extends into the margins. If the top and left margins are zero, the numbering will not be visible. If you still want to see it even at the cost of possible overlap with the text, you will need to define the parameters discussed in the next section. The spacing between vertical and horizontal lines is the same and has a size of $1mm$. An example of such usage can be found in Appendix \ref[implicit].

\sec Defining Parameters
Of course, if it were only possible to create a grid with $1mm$ divisions, it would be very impractical. Therefore, the grid's appearance can be extensively customized using parameters. The first two parameters before the parenthesis specify the size of one cell in the grid. The parameters inside the parenthesis define, in order: left offset, top offset, width, and height of the grid. So the resulting definition looks like this:
\begtt
\grid <horizontal spacing>,<vertical spacing>(<left margin>,<top margin>,
<horizontal size>,<vertical size>)
\endtt
All parameters can be specified in any units supported by \TeX{}.
\begtt
\grid 1cm,10ex(0.5in,2em,400pt,20cm)
\endtt
You can also define parameters using other sequences:
\begtt
\grid \parindent ,\baselineskip(\hoffset ,\voffset ,\hsize ,\vsize)
\endtt
The result of this notation can be seen in Appendix \ref[parametric]. The same result as with the default definition can be obtained by writing:
\begtt
\grid 1mm,1mm(\hoffset ,\voffset ,\hsize ,\vsize)
\endtt
When defining parameters, you can leave some of them empty, in which case their value will be set to the default.

\sec Possible Improvements
I would say the macro has two shortcomings. Both occur only with specific usage.

\secc[background] Background
The implementation uses the token list from \OpTeX{} "\pgbackground". This means that if we wanted to determine the position of something on a background defined by this list, it would not be possible. Depending on the order of definitions, the background will contain either the grid or the original background---never both at the same time.

As a possible improvement, the macro could define its own token list that would be displayed in the same way as "\pgbackground".

\secc Full Page
As soon as we set the left and top offsets to zero, we lose the line numbering. This can be seen in Appendix \ref[whole], where the grid is set to cover the entire page. Thanks to the thick lines, this is not a very serious problem.

One solution would be to display the numbering on all four sides. However, this does not change the visibility of the numbering in example \ref[whole]. Another option would be to display the numbering in small font inside the grid. But this would potentially make it impossible to read exact positions, since the texts would overlap.

\chap Implementation
This chapter describes the process by which I created the macro. I will describe what led me to the decisions I made. First, it is necessary to clarify what the macro should do. Therefore, I will summarize the assumptions I had during development.
\begitems
* easy to use
* configurable parameters
* does not interfere with the text
* clear and readable
\enditems
\sec Interface
Since I was aiming for simplicity, I decided that the grid would be generated by a single command with parameters. I chose 6 of them.

The most important thing is setting the cell size in the grid, since without this feature the grid is useless. The first two parameters handle this. Equally important is the position of the grid's origin, which allows us to easily determine distances from any point on the page. If the origin were fixed, we would have to determine the distance of two points from, say, the edge of the paper and then subtract these distances. This requires two more parameters. The last two are perhaps more of an aesthetic consideration. The grid could always extend to the edge of the paper. For better visibility, however, I chose to make the grid size configurable as well.
The macro definition will therefore look like this:
\begtt
\def\grid#1,#2(#3,#4,#5,#6)
\endtt
\sec Default Parameter Values
Parameters should have some default value for easier handling. I define all parameters for use in auxiliary macros. I will need to numerically compare the height and width to ensure the grid size is respected. Therefore, I define variables for them:
\begtt
\newdimen\h 
\newdimen\w 
\endtt
The others can be defined using "\def".

I detect an empty parameter with the assumption that the parameter will never contain the character "$". This is a good assumption since "$" starts math mode. If the parameter is empty, the condition "\if$#1$" will result in equality after expansion. There I define the default value. The conditions for all parameters look like this:
\begtt
  \if$#1$ \def\x{1mm} \else \def\x{#1} \fi
  \if$#2$ \def\y{1mm} \else \def\y{#2} \fi
  \if$#3$ \def\offsetLeft{\hoffset} \else \def\offsetLeft{#3} \fi
  \if$#4$ \def\offsetTop{\voffset} \else \def\offsetTop{#4} \fi
  \if$#5$ \w=\hsize \else \w=#5 \fi
  \if$#6$ \h=\vsize \else \h=#6 \fi
\endtt
\sec Auxiliary Macros
I define two auxiliary macros for horizontal and vertical lines. Each of them draws a line and calls itself. I want the macro to stop calling itself when it has drawn enough lines with the defined spacing "\x" or "\y" for vertical or horizontal lines respectively, such that another call would exceed the defined width or height of the grid. For this, I need to define variables:
\begtt
\newdimen\curr % filled distance
\newcount\linecount 
\endtt
In the main macro, I call both auxiliary macros once. Before that, I set the variable values to zero.
\begtt
\linecount=0 \curr=0pt\vertical
\linecount=0 \curr=0pt\horizontal
\endtt
\secc Loop
With each call of the auxiliary macro, "\curr" increases by "\x" or "\y" until the size of "\curr" exceeds "\w" or "\h" respectively. I achieve this behavior with a simple condition, which is why "\w" and "\h" had to be defined using "\newdimen":
\begtt
\def\vertical {
  \ifnum\curr<\w
    \advance\curr\x
    \vertical
  \fi
}
\endtt
Inside the loop, I want to achieve different behavior for every tenth line for better visibility during use. I therefore add "\newcount\blockcount", which I also set to zero before calling the auxiliary macro. Every time I draw a thicker, extended tenth line, I increase "\blockcount" by ten. I insert another condition into the previous one. This checks whether "\linecount" equals "\blockcount"---that is when every tenth line is being drawn.
\begtt
  \ifnum\linecount =\blockcount
    \global\advance\blockcount10
    % every 10th line
  \else
    % other lines
  \fi
\endtt
\secc Lines and Numbers
The principle of drawing lines in the correct position is that after drawing each line, I return to the position where I started drawing. Before the line, I output a space of "\curr" as "\hskip" or "\vskip" depending on which auxiliary macro I am in. This ensures correct spacing between lines. The return after output is ensured by zero-dimension "\hbox" and "\vbox" nested inside each other and the commands "\vss" and "\hss" at the end. Printing the number is very similar. Every tenth line is thicker and extended by "\baselineskip", so before calling the "\vertical" or "\horizontal" macro, I calculate the length of this line. For this purpose, I add a new variable "\newdimen\largerline". The first in order is a negative space, which means the number will protrude. The complete "\vertical" will look like this:
\begtt
\def\vertical {\noindent
  \ifnum\curr<\w
    \hbox to0pt{\hskip\offsetLeft\hskip\curr
      \vbox to0pt{\vskip\offsetTop\noindent
      \ifnum\linecount =\blockcount
        \global\advance\blockcount10
        \vbox to0pt{\vss\noindent\hskip2pt\the\linecount\vskip1pt}%
        \vskip-\baselineskip\noindent\vrule height\largerline width 1pt
      \else
        \vrule height\h width 0.1pt
      \fi
    \vss}\hss}%
    \advance\linecount1
    \advance\curr\x
    \ea\vertical
  \fi
}
\endtt
The "\horizontal" macro will look almost identical (see Appendix \ref[code]).

\sec Rendering on Every Page
The current macro could be used standalone on the first page without changing the text, since we return to the starting position after each line. However, "\grid" would then need to be placed at the beginning of every page. But we can only determine the beginning of the page for the first one. Therefore, I decided to use the token list "\_pgbackground" from \OpTeX{}. This comes with the disadvantages mentioned in Section \ref[background]. I therefore add to the main macro:
\begtt
\_pgbackground={\Green
    \largerline=\h \advance\largerline\baselineskip
    \blockcount=0 \linecount=0 \curr=0pt\vertical
    \largerline=\w \advance\largerline\baselineskip
    \blockcount=0 \linecount=0 \curr=0pt\horizontal}%
\endtt
I chose green for the grid color to mimic the appearance of graph paper.
\app Examples of Macro Usage
\ref[implicit] is generated by the command:
\begtt
\grid ,(,,,)
\endtt
\ref[parametric] is generated by the command:
\begtt
\grid \parindent ,\baselineskip(\hoffset ,\voffset ,\hsize ,\vsize)
\endtt
\ref[whole] is generated by the command:
\begtt
\grid 10pt,10pt(0pt,0pt,\pdfpagewidth ,\pdfpageheight)
\endtt
\midinsert
\picw=16cm 
\clabel[implicit]{Example of default grid usage.}
\cinspic example1.pdf
\caption/f Example of default grid usage.
\endinsert
\midinsert
\picw=16cm 
\clabel[parametric]{Example of parameter usage.}
\cinspic example2.pdf
\caption/f Example of parameter usage.
\endinsert
\midinsert
\picw=16cm 
\clabel[whole]{Example of full-page grid usage.}
\cinspic example3.pdf
\caption/f Example of full-page grid usage.
\endinsert

\app[code] Macro Code
\begtt
\newdimen\h 
\newdimen\w 
\newdimen\curr 
\newcount\linecount 
\newcount\blockcount
\newdimen\largerline 
\def\grid#1,#2(#3,#4,#5,#6){%
  \if$#1$ \def\x{1mm} \else \def\x{#1} \fi
  \if$#2$ \def\y{1mm} \else \def\y{#2} \fi
  \if$#3$ \def\offsetLeft{\hoffset} \else \def\offsetLeft{#3} \fi
  \if$#4$ \def\offsetTop{\voffset} \else \def\offsetTop{#4} \fi
  \if$#5$ \w=\hsize \else \w=#5 \fi
  \if$#6$ \h=\vsize \else \h=#6 \fi
  
  \_pgbackground={\Green
    \largerline=\h \advance\largerline\baselineskip
    \blockcount=0 \linecount=0 \curr=0pt\vertical
    \largerline=\w \advance\largerline\baselineskip
    \blockcount=0 \linecount=0 \curr=0pt\horizontal}%
}
\def\vertical {\noindent
  \ifnum\curr<\w
    \hbox to0pt{\hskip\offsetLeft\hskip\curr
      \vbox to0pt{\vskip\offsetTop\noindent
      \ifnum\linecount =\blockcount
        \global\advance\blockcount10
        \vbox to0pt{\vss\noindent\hskip2pt\the\linecount\vskip1pt}%
        \vskip-\baselineskip\noindent\vrule height\largerline width 1pt
      \else
        \vrule height\h width 0.1pt
      \fi
    \vss}\hss}%
    \advance\linecount1
    \advance\curr\x
    \ea\vertical
  \fi
}
\def\horizontal {\noindent
  \ifnum\curr<\h
    \hbox to0pt{\hskip\offsetLeft
      \vbox to0pt{\vskip\offsetTop\vskip\curr
      \ifnum\linecount =\blockcount
        \global\advance\blockcount10
        \hbox to0pt{\hskip-\baselineskip
          \vbox to0pt{\hrule width\largerline height 1pt\vss}\hss}%
        \vbox to0pt{\vskip-\baselineskip\vskip2pt
          \hbox to0pt{\hss\noindent\the\linecount\hskip1pt}\vss}%
      \else
        \hrule width\w height 0.1pt
      \fi
    \vss}\hss}%
    \advance\linecount1
    \advance\curr\y
    \ea\horizontal
  \fi
}
\endtt
\bye